% Figure : règles de Feynman pour la diffusion Raman (fig:feynman_rules)
% Compatible pdflatex. Nécessite \usetikzlibrary{decorations.pathmorphing,arrows.meta}
\begin{figure}[H]
\centering
\begin{tikzpicture}[
    line/.style={thick},
    >=Stealth,
]

% ===================== Colonne gauche : propagateurs =====================

% --- Photon : ligne pointillée ---
\draw[line, dashed] (0,4) -- (2,4);
\node[anchor=west] at (2.3,4) {Photon};

% --- Électron/trou : ligne pleine avec flèche ---
\draw[line, fermion] (0,3) -- (2,3);
\node[anchor=west] at (2.3,3) {Électron};

% --- Phonon : ligne ondulée ---
\draw[line, decorate, decoration={snake, amplitude=1.2mm, segment length=4mm}]
    (0,2) -- (2,2);
\node[anchor=west] at (2.3,2) {Phonon};

% ===================== Colonne droite : vertex =====================

% --- Vertex électron-photon : cercle ---

\draw[fill=white] (7,4) circle (0.2);
\node[anchor=west] at (7.3,4) {Couplage électron-photon $\hat{H}_\text{eem}$};

% --- Vertex électron-défaut : cercle avec croix ---
\draw[fill=white] (7,3) circle (0.2);
\draw[line width=1pt] (6.92,2.92) -- (7.08,3.08);
\draw[line width=1pt] (6.92,3.08) -- (7.08,2.92);
\node[anchor=west] at (7.3,3) {Couplage électron-défaut $\hat{H}_\text{ed}$};

% --- Vertex électron-phonon : carré ---
\def\halfside{0.2} % ajuste cette valeur pour la taille du carré
\draw[fill=white] (7-\halfside,2-\halfside) rectangle (7+\halfside,2+\halfside);
\node[anchor=west] at (7.3,2) {Couplage électron-phonon $\hat{H}_\text{ep}$};

\end{tikzpicture}
\end{figure}